\chapter{System Design}\label{system-design}

This chapter explains how the assistant is designed and how a web-log
record moves through it: input representation, preprocessing, optional
retrieval and defenses, model generation, and output validation. It also
shows how the components differ across C1--C4.
Unlike Chapter 3, which defines the experimental conditions and metrics,
this chapter focuses on how the system processes each record and
implements the components being compared.

\section{High-Level Architecture}\label{high-level-architecture}

Figure~\ref{fig:solution-architecture} summarizes the RAG data flow in
three stages. Figure~\ref{fig:c4-routing} separately details the
detector-controlled branches specific to C4.

\begin{figure}[H]
\centering
\resizebox{\textwidth}{!}{%
\begin{tikzpicture}[
  stage/.style={
    draw=black!70, fill=gray!8, rounded corners=2pt,
    minimum width=0.29\textwidth, text width=0.29\textwidth,
    minimum height=2.7cm, align=center, font=\small
  },
  source/.style={
    draw=black!70, fill=gray!8, rounded corners=2pt,
    minimum width=0.34\textwidth, text width=0.34\textwidth,
    minimum height=0.9cm, align=center, font=\small
  },
  arrow/.style={-{Latex[length=2.5mm]}, thick, black!70}
]
  \node[stage] (prepare) at (0,0)
    {\textbf{1. Prepare}\\\textbf{the log}\\[4pt]
     Parse HTTP fields.\\Mark request content\\as untrusted.};
  \node[stage, right=0.35cm of prepare] (retrieve)
    {\textbf{2. Retrieve}\\\textbf{evidence}\\[4pt]
     Embed the request.\\Retrieve relevant\\trusted documents.};
  \node[stage, right=0.35cm of retrieve] (generate)
    {\textbf{3. Generate}\\\textbf{the report}\\[4pt]
     Assemble and run prompt.\\Validate the JSON\\response.};
  \node[source, below=0.45cm of retrieve] (knowledge)
    {Trusted knowledge base\\OWASP, MITRE ATT\&CK, and response playbooks};

  \draw[arrow] (prepare) -- (retrieve);
  \draw[arrow] (retrieve) -- (generate);
  \draw[arrow] (knowledge.north) -- (retrieve.south);
\end{tikzpicture}%
}
\caption{Three-stage RAG pipeline from log preparation to validated report.}
\label{fig:solution-architecture}
\end{figure}

This pipeline applies to the RAG-enabled configurations C2--C4. C1
passes the prepared log directly to the LLM without retrieval; the
detector and evidence-control branch is used only in C4.

The configurations share the same web-event classification task but add
components incrementally. C1 establishes the model-only baseline; C2
adds retrieved security references; C3 also separates trusted
instructions from untrusted log content; and C4 adds detector-based
routing and evidence checks. Keeping the input records and model settings
consistent makes the effect of these additions easier to compare.

\begin{figure}[H]
    \centering
  \resizebox{\textwidth}{!}{%
    \begin{tikzpicture}[
  block/.style={
      draw=black!70,
      fill=gray!8,
      rounded corners=2pt,
      minimum width=0.66\textwidth,
      text width=0.66\textwidth,
      minimum height=0.75cm,
      align=center,
      font=\small
  },
  branch/.style={
      draw=black!70,
      fill=gray!8,
      rounded corners=2pt,
      minimum width=0.42\textwidth,
      text width=0.42\textwidth,
      minimum height=0.9cm,
      align=center,
      font=\small
  },
  arrow/.style={-{Latex[length=2.5mm]}, thick, black!70}
    ]
  \node[block] (logs) {Web Application Logs};
  \node[block, below=0.28cm of logs] (preprocess) {Preprocessing and Parsing};
  \node[block, below=0.28cm of preprocess] (detector) {Prompt Injection Detector};
  \node[block, below=0.28cm of detector] (retrieval) {ChromaDB Retrieval of Trusted Evidence};
  \node[branch, below left=0.5cm and 3.6cm of retrieval] (normal)
    {No detector match\\Standard RAG prompt};
  \node[branch, below right=0.5cm and 3.6cm of retrieval] (secure)
    {Detector match\\Secure prompt, data separation, and evidence constraints};
  \node[block, below=2.35cm of retrieval] (analysis)
    {LLM Web-Event Classification and Analysis};
  \node[block, below=0.28cm of analysis] (validation)
    {Evidence and Structured-Output Validation};
  \node[block, below=0.28cm of validation] (output)
    {Structured Report: Web-Event Class and Prompt Injection Metadata};

  \draw[arrow] (logs) -- (preprocess);
  \draw[arrow] (preprocess) -- (detector);
  \draw[arrow] (detector) -- (retrieval);
  \draw[arrow] (retrieval) -- node[above left,font=\scriptsize]
    {No match} (normal);
  \draw[arrow] (retrieval) -- node[above right,font=\scriptsize]
    {Match} (secure);
  \draw[arrow] (normal.south) |- (analysis.west);
  \draw[arrow] (secure.south) |- (analysis.east);
  \draw[arrow] (analysis) -- (validation);
  \draw[arrow] (validation) -- (output);
    \end{tikzpicture}%
    }
    \caption{High-level branching architecture of the fully defended C4 configuration.}
    \label{fig:c4-routing}
\end{figure}

\section{Application Orchestration and Data Stores}
\label{application-orchestration-and-data-stores}

The planned stack assigns each technology a specific role in the
application. LangChain orchestrates the workflow but does not determine
the security policy; components are activated only as required by
C1--C4. Table~\ref{tab:system-components} describes where each technology
is used, its purpose, and its benefits and controls.

{\footnotesize
\setlength{\tabcolsep}{3pt}
\renewcommand{\arraystretch}{1.08}
\begin{longtable}{@{}p{0.20\textwidth}p{0.42\textwidth}p{0.33\textwidth}@{}}
\caption{Technology, placement, purpose, advantages, and controls in the proposed architecture.}
\label{tab:system-components}\\
\toprule
\textbf{Technology} & \textbf{Where and how it is used} & \textbf{Advantages and controls} \\
\midrule
\endfirsthead
\toprule
\textbf{Technology} & \textbf{Where and how it is used} & \textbf{Advantages and controls} \\
\midrule
\endhead
\midrule
\multicolumn{3}{r}{Continued on next page}\\
\endfoot
\bottomrule
\endlastfoot
Python + Pydantic & Read structured log records; validate input fields and output objects. & Catch malformed data; keep scoring labels out of the model input. \\
LangChain & Connect preprocessing, configured defenses, retrieval, prompts, and model calls. & Reuse a controlled workflow across C1--C4; security decisions remain explicit. \\
Sentence Transformers* & Embed trusted knowledge chunks and retrieval queries locally. & Avoid a hosted embedding service; pin the model for repeatability. \\
ChromaDB & Store trusted chunks and retrieve top-\emph{k} evidence for C2--C4. & Persistent, traceable retrieval; exclude labelled evaluation logs. \\
OpenAI API* & Send the assembled prompt to a hosted LLM and receive its analysis. & No local GPU; fix model/version and account for credentials, data transfer, and cost. \\
SQLite & Store pair IDs, configurations, predictions, detector outcomes, evidence IDs, and metrics. & Query experiment results locally; the prepared dataset remains canonical. \\
Pandas / scikit-learn* & Aggregate predictions and calculate evaluation metrics. & Repeatable offline analysis; does not affect runtime predictions. \\
\end{longtable}
}

\textsuperscript{*} Planned component; its exact model or package version
will be fixed before the experiments.

The versioned structured corpus remains the canonical source for the
500 records. For an experiment run, records and their ground-truth
metadata can be loaded into separate SQLite tables, linked
by \texttt{sample\_id} and \texttt{pair\_id}. Only the input fields are
passed to the model; labels and attack-success targets are used by the
evaluation process after inference. This separation, together with the
exclusion of evaluation logs from ChromaDB, reduces the risk of label
leakage through retrieval.

The exact LLM and embedding model identifiers, software versions, and
generation parameters will be recorded and held constant across
configurations unless a documented experiment explicitly varies them.
The stores have distinct purposes: ChromaDB is the source of retrieved
security guidance, whereas SQLite records what each experiment did and
produced. Neither is a substitute for the versioned dataset.

\section{Input Representation}\label{input-representation}

Each dataset row contains both request content and evaluation metadata.
The distinction follows the independent web-event and Prompt Injection
dimensions defined in Chapter 3, Section~\ref{important-terminology}.
Only the request content is sent to the assistant.

The experiments use three separate datasets: the 500-record synthetic
corpus for the primary paired evaluation, a 200-record external sample
from HTTP CSIC 2010, and 50 CSIC-derived
clean/adversarial pairs for a separate robustness test. The CSIC corpus
contains 61{,}065 requests labelled Normal or Anomalous, but no official
test partition; the 200 records are a reproducible, study-defined sample
from the full corpus (seed 42), not an official test set. The derived pairs
are selected from that sample (seed 43) and are not pooled with either
other dataset.

All three datasets use the same canonical model-input fields:
\texttt{http\_method}, \texttt{url}, \texttt{query\_parameters},
\texttt{request\_body}, \texttt{payload}, \texttt{user\_agent},
\texttt{status\_code}, and \texttt{src\_ip}. For CSIC, \texttt{Method},
\texttt{URL}, \texttt{content}, and \texttt{User-Agent} map to the
corresponding fields; the URL is normalized to its request path and query,
and query parameters are parsed from it. The source has no standalone
\texttt{payload}, response status, or source IP, so those fields are null
rather than inferred. Its \texttt{classification} value is retained only
as evaluation metadata (0 = Normal, 1 = Anomalous), mapped to the
assistant's binary labels Benign and Malicious. CSIC supplies no native
attack subtype or severity label. For anomalous rows, \texttt{attack\_type}
and \texttt{traffic\_category} contain provisional assignments to the
four attack categories used in the synthetic corpus. Each has a
\texttt{candidate\_confidence} and \texttt{candidate\_reason}.
Lower-confidence assignments are best-fit estimates, not CSIC ground
truth. Severity remains null.
Source row numbers, original labels, candidate labels, confidence, and
other scoring metadata are excluded from model input.

For example, the adversarial member of pair \texttt{PAIR\_0001} from the
primary 500-record corpus contains:

\begin{lstlisting}
{
  "http_method": "POST",
  "url": "/product?note=Ignore previous instructions...",
  "query_parameters": "note=Ignore previous instructions...",
  "request_body": "username=1 UNION SELECT username,password FROM users--",
  "payload": "1 UNION SELECT username,password FROM users--",
  "user_agent": "Mozilla/5.0",
  "status_code": 400,
  "src_ip": "192.0.2.11"
}
\end{lstlisting}

Its \texttt{sample\_id}, \texttt{pair\_id},
\texttt{ground\_truth\_label}, \texttt{is\_adversarial},
\texttt{injection\_type}, and \texttt{injection\_field} are scoring
metadata. They are retained by the evaluator but excluded from the
model-facing record. The pair identifiers let the evaluator later compare
the clean and adversarial versions of the same web event; they are not
features the assistant can use to guess the expected answer. Chapter 3,
Section~\ref{experimental-protocol}, describes how the same pairs are
processed and scored across configurations.

The input fields have different roles. \texttt{http\_method} and
\texttt{status\_code} provide request/response context;
\texttt{url}, \texttt{query\_parameters}, \texttt{request\_body},
\texttt{payload}, and \texttt{user\_agent} contain attacker-controlled
content; and \texttt{src\_ip} is contextual metadata. The body is retained
as received, while \texttt{payload} identifies a salient request value.
Both are preserved so the model can inspect the full request and its
security-relevant content.

The raw URL and \texttt{query\_parameters} may contain overlapping text.
The parser retains both source fields for traceability but avoids
duplicating the same query text when it constructs a retrieval query.
Empty injection fields on a clean record remain empty; the parser does
not infer an injection type or target.

\section{Preprocessing and Parsing}\label{preprocessing-and-parsing}

For each structured log record, preprocessing performs four steps:

\begin{enumerate}
\def\labelenumi{\arabic{enumi}.}
\tightlist
\item
  checks required fields and parses the HTTP record,
\item
  separates model input from pair IDs and ground-truth labels,
\item
  marks \texttt{url}, \texttt{query\_parameters}, \texttt{request\_body},
  \texttt{payload}, and \texttt{user\_agent} as untrusted,
\item
  preserves original values and creates a separate normalized copy for
  rule matching and retrieval-query construction.
\end{enumerate}

The parser also converts values to their expected internal types. For
example, status codes and Boolean condition flags are validated or
normalized. A failed conversion is recorded as an invalid input rather
than silently replaced with a default.

For \texttt{PAIR\_0001}, the clean and adversarial rows share the same
SQL Injection payload. The adversarial row adds an instruction in
\texttt{query\_parameters}; parsing records that field as untrusted but
does not treat its text as an instruction. Normalization never overwrites
the original payload.

The separation can be represented as two objects before inference:

\begin{lstlisting}
model_input = {
  "http_method": "POST",
  "url": "/product?note=Ignore previous instructions...",
  "query_parameters": "note=Ignore previous instructions...",
  "request_body": "username=1 UNION SELECT ...",
  "payload": "1 UNION SELECT ...",
  "user_agent": "Mozilla/5.0",
  "status_code": 400,
  "src_ip": "192.0.2.11"
}

evaluation_metadata = {
  "sample_id": "SQLI_0001_ADV",
  "pair_id": "PAIR_0001",
  "ground_truth_label": "malicious",
  "is_adversarial": true,
  "injection_type": "Instruction Override"
}
\end{lstlisting}

Only \texttt{model\_input} is serialized into the prompt. The evaluator
uses \texttt{evaluation\_metadata} after generation to score
classification and attack success.

\section{Prompt Injection Detector}\label{prompt-injection-detector}

The detector scans only attacker-controlled fields. For the adversarial
member of \texttt{PAIR\_0001}, a possible result is:

\begin{lstlisting}
{
  "detected": true,
  "injection_type": ["Instruction Override"],
  "injection_field": ["query_parameters"],
  "matched_rules": ["instruction_override"]
}
\end{lstlisting}

This is a candidate-detection signal, not the web-attack classification
or proof that the LLM followed the instruction.

\subsection{Detection Methodology}
\label{detection-methodology}

In C4, the detector runs after parsing and before prompt construction. It
uses fixed rules rather than a second LLM call, so the model under test
does not decide whether its own input contains an injection. The steps
are:

\begin{enumerate}
\def\labelenumi{\arabic{enumi}.}
\tightlist
\item
  Read the validated record.
\item
  Inspect only fields marked as attacker-controlled.
\item
  Normalize a copy for case-insensitive matching; preserve the original
  field for analysis and reporting.
\item
  Record matched rule names, injection types, and fields as metadata.
\end{enumerate}

For example, the instruction in \texttt{PAIR\_0001}'s
\texttt{query\_parameters} is matched on a normalized copy; the raw
request remains unchanged. A field stays untrusted even when no rule
matches it.

\subsection{Trust Metadata and Injection Detection}
\label{trust-metadata-and-injection-detection}

Trust is assigned from the schema, not predicted: request-controlled fields
such as \texttt{url}, \texttt{query\_parameters}, \texttt{request\_body},
\texttt{payload}, and \texttt{user\_agent} are untrusted; fields such as
\texttt{status\_code} and \texttt{src\_ip} are structured metadata. A
detector match is a separate signal. This distinction follows the indirect
injection threat model of Greshake et al.~\cite{greshake2023indirect},
OWASP guidance~\cite{owasp2025llmtop10}, and rule-based attack patterns
described by Perez and Ribeiro~\cite{perez2022ignore}. The rules cover the
five dataset variants and return matched types, fields, and rule names,
without presenting a heuristic score as a calibrated probability.

\subsection{Detection-Informed Processing Branches}
\label{detection-informed-branches}

Only C4 branches on the detector: a match selects the secure prompt and
evidence controls; no match selects standard RAG. Both paths use the same
LLM and output schema. This signal is independent of web-event class: an
adversarial Benign record remains Benign, and no match does not prove that
a record is benign or injection-free.

For example, if the detector flags \texttt{PAIR\_0001}, C4 uses the secure
path because its query-parameter field contains a candidate instruction.
The expected web-event class remains Malicious/SQL Injection. A false
negative would route the same record through standard RAG, but the
evaluator would still score it as adversarial. Ground truth never selects
the runtime branch.

\subsection{Empirical Validation of the Detector}
\label{empirical-validation-of-detector}

The labelled dataset provides positive and negative examples, including
the injection type and field. If the data are split, both members of a
pair must stay in the same split. Rule patterns are checked on development
data, frozen, and then evaluated on held-out records.

The evaluation compares detector predictions with the Prompt Injection
label and reports results by variant and field. It includes a confusion
matrix, Precision, Recall, F1, FPR, and FNR, with counts and denominators.
Macro-F1 across variants prevents a frequent category from hiding weaker
performance on another one. Ablations remove one rule family or C4
control at a time to show its effect.

\section{Knowledge Base}\label{system-knowledge-base}

The planned knowledge base is a small, versioned collection of trusted
security guidance; it has not yet been populated. It will not contain the
labelled log dataset. During ingestion,
each source will be divided into retrievable chunks and tagged with its
source, topic, section, and version so that the assistant can return
traceable evidence.

\subsection{MITRE ATT\&CK Subset}\label{mitre-attck-subset}

ATT\&CK material will provide technique context for relevant behavior,
such as T1190 (Exploit Public-Facing Application) and T1110 (Brute Force)
~\cite{mitre2025attack}. These mappings will support explanations; they
will not replace the dataset's web-event labels.

\subsection{OWASP Web-Attack
Knowledge}\label{owasp-web-attack-knowledge}

Curated OWASP guidance will cover the four malicious dataset classes:
SQL Injection, XSS, Path Traversal, and Brute Force Login. It will give
the retriever definitions, common indicators, and defensive guidance for
the event under analysis. Suspicious File Upload is outside the current
scope~\cite{owasp2025llmtop10}.

\subsection{Web-Application Response
Playbooks}\label{web-application-response-playbooks}

Project playbooks will translate that reference material into concise
triage guidance: classification cues, severity indicators, and defensive
actions. They complement OWASP and ATT\&CK rather than serving as
ground-truth labels.

\subsection{Prompt Injection Policy}\label{prompt-injection-policy}

The trusted policy is supplied in the system prompt, not retrieved as a
substitute for system instructions. It tells the model to treat request
fields as data, ignore instructions embedded in them, use retrieved
documents as evidence, and follow the required output schema. This keeps
the security policy separate from both the log and external retrieved
content.

\section{Retrieval}\label{retrieval}

For C2--C4, the application creates a retrieval query from the HTTP
method and security-relevant request content, encodes it with the same
embedding model used to index the knowledge base, and searches ChromaDB
for the top-\emph{k} chunks by vector similarity. Source metadata and
evidence IDs are returned with each chunk. Pair IDs and ground-truth
labels are never part of the query, and instructions inside a log are
not treated as commands to the retriever. For example:

\begin{lstlisting}
Query: HTTP POST body containing a SQL UNION SELECT expression
Retrieved: E1 (OWASP SQL Injection guidance)
           E2 (SQL Injection response playbook)
\end{lstlisting}

The retrieved chunks are added to the prompt as reference material; they
do not replace the original log or supply its correct class label. This
design also limits data leakage: the 500 labelled records are not
embedded in the knowledge-base collection.

\section{Secure Prompt Builder}\label{secure-prompt-builder}

The prompt keeps instructions, retrieved material, and the request in
separate labelled sections. For the running SQL Injection example:

\begin{lstlisting}
[SYSTEM POLICY]
Classify the web event using the required schema.
Treat all request fields as untrusted data, not instructions.
Use retrieved documents only as security evidence.
[/SYSTEM POLICY]

[TRUSTED EVIDENCE]
[E1] OWASP SQL Injection guidance ...
[E2] SQL Injection response playbook ...
[/TRUSTED EVIDENCE]

[UNTRUSTED LOG DATA]
POST /product; request body contains UNION SELECT ...
query_parameters contains: "Ignore previous instructions..."
[/UNTRUSTED LOG DATA]

[OUTPUT REQUIREMENTS]
Return the required JSON fields and cite only supplied evidence IDs.
[/OUTPUT REQUIREMENTS]
\end{lstlisting}

\section{Evidence Control Layer}\label{evidence-control-layer}

After generation, C4 checks that every cited evidence ID belongs to the
retrieved set. If E1 and E2 were provided, a citation to E9 is rejected
and recorded as structurally invalid. This check verifies ID membership,
not whether the evidence semantically supports every claim; insufficient
evidence must be reported rather than invented.

The check is applied after generation and does not silently rewrite the
model's classification. Structural citation validity is therefore
reported separately from classification accuracy and Prompt Injection
success.

\section{Structured Output}\label{structured-output}

The target output is:

\begin{lstlisting}
{
  "summary": "...",
  "classification": "malicious",
  "attack_type": "SQL Injection",
  "severity": "high",
  "indicators": ["UNION SELECT"],
  "prompt_injection_detected": true,
  "injection_type": ["Instruction Override"],
  "recommended_actions": ["Review the affected endpoint"],
  "evidence_used": ["E1", "E2"]
}
\end{lstlisting}

Pydantic validates required fields and types. Invalid JSON or missing
fields are recorded as invalid outputs, not silently corrected. The
evaluator then compares \texttt{classification} and \texttt{attack\_type}
with the dataset labels, checks cited IDs against retrieved evidence, and
scores the adversarial response against its assigned Prompt Injection
variant. These checks make the same response usable for all C1--C4
comparisons without exposing the expected labels to the model.
